\section{Resumen ejecutivo}

Lo que se quiere conseguir a la hora de realizar este trabajo, es conseguir reunir en una única página web los lugares de nuestro planeta en los que se han conseguido certificar o validar avistamientos de ovnis.Mediante la utilización de una amplia gama de fuentes de datos heterogéneos pertenecientes a diferentes satélites y empresas se quiere lograr trazar en un mapa global en tres dimensiones dichos datos de tal forma que se pueda ver la ubicación del avistamiento y los datos relevantes a ese suceso.

Estas amplias fuentes de datos cuentan con datos en diferentes formatos, por lo que se cuenta con datos auditivos, imágenes, video y demás. Además se quiere utilizar los datos pertenecientes a distintas compañías, tanto públicas como privadas con el fin de contrastar la información obtenida mediante diferentes fuentes sobre cada avistamiento.

Además, gracias al estado actual de la inteligencia artificial, se puede hacer uso de la misma con el fin de agilizar el trabajo ya que por ejemplo, si contamos con audios muy largos podemos utilizar esta tecnología para que seleccione las secciones más interesantes de ese audio.
